\scanpage{178}
\mspage{13}
% archive sheet number: 177
Posons $I'=\pi_1(U',\xi')$ resp. $I'=\pi_1(\eta,\xi')$. On a (SGA 1 \uline{XIII})
\[
\pi_1^{\unsure{(p')|(p)}}(X'_{\bar\eta'},\xi')\mathrel{\overset{\sim}{\text{---}}}\pi_1^{(p')}(X_{\bar\eta},\xi)
\]
\deleted{\unsure{donc}} et d'autre part, $U'\mathrel{\text{---}}U$ (resp. $\eta'\mathrel{\text{---}}\eta$) étant connexes, on a
\[
I'\unsure{\mathrel{\text{---}}|\longrightarrow}I\qquad\text{\uline{surjectif}}.
\]
Donc il suffit de prouver que l'action \deleted{extérieure} de $I'$ sur $\pi_1^{(p')}(X'_{\bar\eta'},\xi')$ \deleted{\unsure{a les vertus dites}, \unsure{ou}} \inserted{a les vertus annoncées}.
\deleted{\unsure{ou le cas où on} dispose une section $g$ de $\bar X/S$ dans l'aire de lissité, qui définit une section de $X'$ sur $S'$ dans l'aire de lissité, pour prouver l'assertion faite sur l'action de $I$ sur $\pi_1^{(p')}(X_{\bar\eta},\xi)$; dans ce dernier cas, il \deleted{\unsure{sera}} faudra modifier la construction précédente en prenant $H$ passant par la section, et générique (et ``générique'' pour \deleted{\unsure{assurer}} \inserted{autres propriétés}). D'où D'autre part, comme les composantes irréductibles de}

D'autre part, comme $X_{\bar\eta}^{\deleted{\unsure{(1)|*}}}$ est irréductible de dim $\geq2$, et $X^{(1)}_{\bar\eta}$ en est une section hyperplane \uline{générique}, le th. de Bertini \inserted{et le th. d'invariance de $\pi_1^{(p')}$ \deleted{\unsure{géom}} par ext. du corps de base} impliquent que
\[
\pi_1^{(p')}(X^{(1)}_{\bar\eta},\xi)\mathrel{\text{---}}\pi_1^{(p')}(X_{\bar\eta},\xi)
\]
est \uline{surjectif}. Donc, on est ramené à prouver le même énoncé 1.4 pour $\bar X^{(1)},Y^{(1)},U'$ resp. $\eta'$.
Cela achève la réduction au cas de la dim 1.
